% Written last, from the results. 150-250 words, no citations, no undefined acronyms.
Forecasts of when a price trend will end are widely used by investors, but their skill is rarely
compared with what the definition of a trend implies by itself. We forecast the end of price swings
of Vietnamese stocks, defined by a directional-change rule whose threshold scales with each stock's
volatility, using landmark survival models and features timed to what was public before the market
opened. Models were designed on 2012--2016 and tested once, on 2017--2026, for the members of the VN30
index and all listed banks. Reversal timing is predictable out of sample: the main model reaches a
concordance index of \CABCDx{} against 0.5 for the Kaplan--Meier benchmark and \CRW{} for a
random-walk null, with a lower integrated Brier score. Most of this skill is mechanical. The distance to the reversal
threshold alone accounts for \MechShareDIST\% of the gain, and a flexible model of swing age and
distance is as well calibrated as the richer linear models. In the pre-registered linear test,
stock and local market features improve the ranking of swings by \GainOverMFlex{} in the
concordance index but not their predicted probabilities; post hoc, non-linear models of the same
features improve both, by small margins. Overnight global markets, bank fundamentals and news bring
no detectable improvement, nor does swing age once the distance is known, and an exit rule based on the forecasts earns a lower Sharpe
ratio than buy-and-hold after costs. Turning-point forecasts should be judged against a mechanical
benchmark built from the labelling rule.
